\documentclass[UTF8,11pt,a4paper]{ctexart}

\usepackage[a4paper,margin=2.2cm]{geometry}
\usepackage{amsmath,amssymb,mathtools,mathrsfs}
\usepackage{enumitem}
\usepackage[dvipsnames]{xcolor}
\usepackage[most]{tcolorbox}
\usepackage{fancyhdr}

\pagestyle{fancy}
\fancyhf{}
\lhead{信息理论第一次作业}
\cfoot{\thepage}
\setlength{\headheight}{14pt}

\setlist[enumerate]{leftmargin=2.2em}

\newtcolorbox{problem}[1]{
  enhanced,
  breakable,
  colback=white,
  colframe=MidnightBlue!75,
  colbacktitle=MidnightBlue,
  coltitle=white,
  boxrule=0.7pt,
  arc=1.5mm,
  left=5mm,right=5mm,top=3mm,bottom=3mm,
  title=\textbf{#1}
}

\begin{document}

\begin{center}
  {\zihao{3}\bfseries 信息理论第一次作业}\\[0.4em]
  {\small 以下各题对数均取 $\log_2$}
\end{center}

\vspace{0.8em}

\begin{problem}{2.1}
A村有一半人说真话，$\frac{3}{10}$ 人总说假话，$\frac{2}{10}$ 人拒绝回答；B村有 $\frac{3}{10}$ 人诚实，一半人说谎，$\frac{2}{10}$ 人拒绝回答。现随机地从A村和B村抽取人，$p$ 为抽到A村人的概率，$1-p$ 为抽到B村人的概率。求通过测试某人说话状态所得到的关于其所属村庄的平均信息量，并求其最大值。
\end{problem}
\noindent\textbf{解：}
设村庄变量为 $V\in\{A,B\}$，测试结果为 $S\in\{\text{真},\text{假},\text{拒}\}$。则
\[
P(S|A)=\left(\frac12,\frac{3}{10},\frac15\right),\qquad
P(S|B)=\left(\frac{3}{10},\frac12,\frac15\right).
\]
所以
\[
P(S)=\left(0.3+0.2p,\ 0.5-0.2p,\ 0.2\right).
\]
平均得到的信息量即互信息
\[
I(V;S)=H(S)-H(S|V).
\]
而
\[
H(S|V)=H\!\left(\frac12,\frac{3}{10},\frac15\right),
\]
因为两村对应分布只是前两项互换，熵相同。故
\[
I(V;S)=H(0.3+0.2p,\ 0.5-0.2p,\ 0.2)-H\!\left(\frac12,\frac{3}{10},\frac15\right).
\]
当第三项固定为 $0.2$ 时，熵在前两项相等时取最大，即
\[
0.3+0.2p=0.5-0.2p \quad\Longrightarrow\quad p=\frac12.
\]
于是
\[
I_{\max}=H(0.4,0.4,0.2)-H(0.5,0.3,0.2)\approx 0.03645\ \text{bit}.
\]
\vspace{0.8em}

\begin{problem}{2.2}
一个无偏骰子掷一次：若出现 $1,2,3,4$ 点，则投一枚均匀硬币一次；若出现 $5,6$ 点，则投同一硬币两次。求硬币投掷中正面出现次数对于骰子点数所提供的信息。
\end{problem}
\noindent\textbf{解：}
设骰子结果为 $D$，正面出现次数为 $K$。则
\[
P(K|D=1,2,3,4)=\left(\frac12,\frac12,0\right),
\qquad
P(K|D=5,6)=\left(\frac14,\frac12,\frac14\right),
\]
这里 $K=0,1,2$。

先求
\[
P(K=0)=\frac46\cdot\frac12+\frac26\cdot\frac14=\frac{5}{12},
\]
\[
P(K=1)=\frac46\cdot\frac12+\frac26\cdot\frac12=\frac12,
\qquad
P(K=2)=\frac26\cdot\frac14=\frac{1}{12}.
\]
故
\[
H(K)=H\!\left(\frac{5}{12},\frac12,\frac{1}{12}\right).
\]
又
\[
H(K|D)=\frac46\cdot 1+\frac26\cdot\frac32=\frac76.
\]
因此
\[
I(D;K)=H(K)-H(K|D)
=H\!\left(\frac{5}{12},\frac12,\frac{1}{12}\right)-\frac76
\approx 0.15834\ \text{bit}.
\]
\vspace{0.8em}

\begin{problem}{2.4}
随机掷 $3$ 颗骰子，以 $X$ 表示第一颗骰子的结果，以 $Y$ 表示前两颗骰子点数之和，以 $Z$ 表示三颗骰子点数之和。求 $H(X|Y)$，$H(Y|X)$，$H(Z|X,Y)$，$H(X,Z|Y)$ 和 $H(Z|X)$。
\end{problem}
\noindent\textbf{解：}
记三颗骰子分别为 $D_1,D_2,D_3$，则
\[
X=D_1,\qquad Y=D_1+D_2,\qquad Z=D_1+D_2+D_3.
\]

\[
H(Y|X)=H(D_2)=\log 6.
\]

\[
H(Z|X,Y)=H(D_3)=\log 6.
\]

对 $y=2,3,\dots,12$，给定 $Y=y$ 时，$X$ 的可能个数为
\[
n(y)=1,2,3,4,5,6,5,4,3,2,1.
\]
且在这些可能值上等概，故
\[
H(X|Y)=\sum_{y=2}^{12}\frac{n(y)}{36}\log n(y)
=\frac{2(2\log 2+3\log 3+4\log 4+5\log 5)+6\log 6}{36}
\approx 1.89552.
\]

由链式法则，
\[
H(X,Z|Y)=H(X|Y)+H(Z|X,Y)\approx 1.89552+\log 6\approx 4.48049.
\]

再由 $Z|X$ 等价于 $D_2+D_3$ 的分布，
\[
P(D_2+D_3=s)=\frac{n(s)}{36},\qquad s=2,\dots,12.
\]
因此
\[
H(Z|X)=H(D_2+D_3)=\log 36-H(X|Y)\approx 3.27440.
\]

综上，
\[
H(X|Y)\approx 1.89552,\quad H(Y|X)=\log 6\approx 2.58496,
\]
\[
H(Z|X,Y)=\log 6\approx 2.58496,\quad H(X,Z|Y)\approx 4.48049,
\]
\[
H(Z|X)\approx 3.27440.
\]
\vspace{0.8em}

\begin{problem}{2.5}
设一个系统传送 $10$ 个数字 $0,1,2,\dots,9$。奇数在传送时以 $0.5$ 概率等可能地错成另外的奇数，而其他数字总能正确接收。求收到一个数字后平均得到的信息量。
\end{problem}
\noindent\textbf{解：}
设发送数字为 $X$，接收数字为 $Y$，并设 $X$ 在 $0,\dots,9$ 上等概。

偶数一定正确传送；若发送奇数，则输出分布为
\[
\left(\frac12,\frac18,\frac18,\frac18,\frac18\right),
\]
故其条件熵为
\[
-\frac12\log\frac12-4\cdot\frac18\log\frac18=2.
\]
于是
\[
H(Y|X)=\frac12\cdot 0+\frac12\cdot 2=1.
\]

再看输出分布：每个偶数以概率 $0.1$ 收到；每个奇数也有
\[
0.1\cdot \frac12+4\cdot 0.1\cdot \frac18=0.1,
\]
所以 $Y$ 在 $0,\dots,9$ 上仍等概，从而
\[
H(Y)=\log 10.
\]
平均得到的信息量为
\[
I(X;Y)=H(Y)-H(Y|X)=\log 10-1\approx 2.32193\ \text{bit}.
\]
\vspace{0.8em}

\begin{problem}{2.6}
对任意概率分布的随机变量，证明：
\begin{enumerate}[label=\textnormal{(\alph*)}]
  \item $H(X|Y)+H(Y|Z)\ge H(X|Z)$；
  \item $\dfrac{H(X|Y)}{H(X,Y)}+\dfrac{H(Y|Z)}{H(Y,Z)}\ge\dfrac{H(X|Z)}{H(X,Z)}$。
\end{enumerate}
\end{problem}
\noindent\textbf{解：}
\begin{enumerate}[label=\textnormal{(\alph*)}]
  \item
  \[
  H(X|Z)\le H(X,Y|Z)=H(Y|Z)+H(X|Y,Z)\le H(Y|Z)+H(X|Y).
  \]
  故
  \[
  H(X|Y)+H(Y|Z)\ge H(X|Z).
  \]

  \item 记
  \[
  A=H(X|Y),\qquad B=H(Y|Z),\qquad C=H(Z).
  \]
  则
  \[
  H(X,Y)=A+H(Y)\le A+B+C,\qquad H(Y,Z)=B+C\le A+B+C.
  \]
  因而
  \[
  \frac{H(X|Y)}{H(X,Y)}+\frac{H(Y|Z)}{H(Y,Z)}
  \ge
  \frac{A}{A+B+C}+\frac{B}{A+B+C}
  =
  \frac{A+B}{A+B+C}.
  \]
  由 (a) 知 $A+B\ge H(X|Z)$，又函数 $f(x)=\dfrac{x}{x+C}$ 单调递增，所以
  \[
  \frac{A+B}{A+B+C}\ge \frac{H(X|Z)}{H(X|Z)+C}
  =\frac{H(X|Z)}{H(X,Z)}.
  \]
  结论成立。
\end{enumerate}
\vspace{0.8em}

\begin{problem}{2.9}
若 $3$ 个随机变量 $X,Y,Z$ 满足 $X+Y=Z$，其中 $X$ 和 $Y$ 独立，试证
\begin{enumerate}[label=\textnormal{(\alph*)}]
  \item $H(X)\le H(Z)$；
  \item $H(Y)\le H(Z)$；
  \item $H(X,Y)\ge H(Z)$；
  \item $I(X;Z)=H(Z)-H(Y)$；
  \item $I(X,Y;Z)=H(Z)$；
  \item $I(X;Y,Z)=H(X)$；
  \item $I(Y;Z|X)=H(Y)$；
  \item $I(X;Y|Z)=H(X|Z)=H(Y|Z)$。
\end{enumerate}
\end{problem}
\noindent\textbf{解：}
由 $Z=X+Y$ 可知：给定 $X$ 与 $Z$ 可唯一确定 $Y$，给定 $Y$ 与 $Z$ 可唯一确定 $X$。

\begin{enumerate}[label=\textnormal{(\alph*)}]
  \item 由
  \[
  H(Z|Y)=H(X|Y)=H(X)
  \]
  及 $H(Z)=H(Z|Y)+I(Y;Z)$，得
  \[
  H(X)\le H(Z).
  \]

  \item 同理，
  \[
  H(Z|X)=H(Y|X)=H(Y),
  \]
  故
  \[
  H(Y)\le H(Z).
  \]

  \item 由 (e) 将证明
  \[
  H(Z)=I(X,Y;Z)\le H(X,Y),
  \]
  即
  \[
  H(X,Y)\ge H(Z).
  \]

  \item
  \[
  I(X;Z)=H(Z)-H(Z|X)=H(Z)-H(Y).
  \]

  \item 因为 $Z$ 是 $(X,Y)$ 的确定函数，
  \[
  H(Z|X,Y)=0,
  \]
  所以
  \[
  I(X,Y;Z)=H(Z)-H(Z|X,Y)=H(Z).
  \]

  \item 由 $X=Z-Y$，
  \[
  H(X|Y,Z)=0,
  \]
  故
  \[
  I(X;Y,Z)=H(X)-H(X|Y,Z)=H(X).
  \]

  \item
  \[
  I(Y;Z|X)=H(Y|X)-H(Y|X,Z)=H(Y)-0=H(Y).
  \]

  \item
  \[
  I(X;Y|Z)=H(X|Z)-H(X|Y,Z)=H(X|Z).
  \]
  又因为 $(X,Z)$ 与 $(Y,Z)$ 一一对应，
  \[
  H(X,Z)=H(Y,Z),
  \]
  从而
  \[
  H(X|Z)=H(Y|Z).
  \]
  所以
  \[
  I(X;Y|Z)=H(X|Z)=H(Y|Z).
  \]
\end{enumerate}
\vspace{0.8em}

\begin{problem}{2.12}
证明下列关系：
\begin{enumerate}[label=\textnormal{(\alph*)}]
  \item $H(Y,Z|X)\le H(Y|X)+H(Z|X)$，等号成立的充要条件为对一切满足 $p(x_i)>0$ 的 $i$ 及任意 $j,k$，有
  \[
  p(y_j,z_k|x_i)=p(y_j|x_i)p(z_k|x_i);
  \]
  \item $H(Y,Z|X)=H(Y|X)+H(Z|X,Y)$；
  \item $H(Z|X,Y)\le H(Z|Y)$，等号成立的充要条件为对一切满足 $p(y_j)>0$ 的 $j$ 及任意 $i,k$，有
  \[
  p(x_i,z_k|y_j)=p(x_i|y_j)p(z_k|y_j).
  \]
\end{enumerate}
\end{problem}
\noindent\textbf{解：}
\begin{enumerate}[label=\textnormal{(\alph*)}]
  \item 由条件熵链式法则，
  \[
  H(Y,Z|X)=H(Y|X)+H(Z|X,Y).
  \]
  又因为增加条件不会增大条件熵，即
  \[
  H(Z|X,Y)\le H(Z|X),
  \]
  所以
  \[
  H(Y,Z|X)\le H(Y|X)+H(Z|X).
  \]

  并且
  \[
  H(Y|X)+H(Z|X)-H(Y,Z|X)=I(Y;Z|X)\ge 0.
  \]
  等号成立当且仅当 $I(Y;Z|X)=0$，也就是给定 $X$ 时 $Y,Z$ 条件独立，即对一切满足 $p(x_i)>0$ 的 $i$ 及任意 $j,k$，
  \[
  p(y_j,z_k|x_i)=p(y_j|x_i)p(z_k|x_i).
  \]

  \item 由条件熵链式法则，
  \[
  H(Y,Z|X)=H(Y|X)+H(Z|X,Y).
  \]

  \item
  \[
  H(Z|Y)-H(Z|X,Y)=I(X;Z|Y)\ge 0,
  \]
  所以
  \[
  H(Z|X,Y)\le H(Z|Y).
  \]
  等号成立当且仅当 $I(X;Z|Y)=0$，即给定 $Y$ 时 $X,Z$ 条件独立：
  \[
  p(x_i,z_k|y_j)=p(x_i|y_j)p(z_k|y_j).
  \]
\end{enumerate}
\vspace{0.8em}

\begin{problem}{2.13}
令 $X_1$ 和 $X_2$ 为分别定义在字符表
\[
\mathscr{X}_1=\{1,2,\dots,m\},\qquad
\mathscr{X}_2=\{m+1,\dots,n\}
\]
上的两个离散随机变量，它们的分布函数为 $p_1(\cdot)$ 和 $p_2(\cdot)$。现构造随机变量
\[
X=
\begin{cases}
X_1, & \text{以概率 }\alpha,\\
X_2, & \text{以概率 }1-\alpha.
\end{cases}
\]
\begin{enumerate}[label=\textnormal{(\alph*)}]
  \item 试用 $H(X_1)$、$H(X_2)$ 和 $\alpha$ 来表示 $X$ 的熵 $H(X)$。
  \item 在 $\alpha$ 上求 $H(X)$ 的极大值，证明
\[
2^{H(X)}\le 2^{H(X_1)}+2^{H(X_2)}.
\]
\end{enumerate}
\end{problem}
\noindent\textbf{解：}
设选择变量
\[
U=
\begin{cases}
1, & \text{选 }X_1,\\
2, & \text{选 }X_2.
\end{cases}
\]
则 $P(U=1)=\alpha$，$P(U=2)=1-\alpha$。由于 $\mathscr{X}_1,\mathscr{X}_2$ 不相交，所以由 $X$ 可唯一确定 $U$，故
\[
H(X)=H(U)+H(X|U).
\]
于是
\[
H(X)=h(\alpha)+\alpha H(X_1)+(1-\alpha)H(X_2),
\]
其中
\[
h(\alpha)=-\alpha\log\alpha-(1-\alpha)\log(1-\alpha).
\]

对 $\alpha$ 求导：
\[
\frac{\mathrm{d}H(X)}{\mathrm{d}\alpha}
=\log\frac{1-\alpha}{\alpha}+H(X_1)-H(X_2).
\]
令导数为 $0$，得
\[
\alpha^\ast=\frac{2^{H(X_1)}}{2^{H(X_1)}+2^{H(X_2)}}.
\]
此时取得极大值，且
\[
H_{\max}(X)=\log\!\bigl(2^{H(X_1)}+2^{H(X_2)}\bigr).
\]
因此对任意 $\alpha$，
\[
H(X)\le \log\!\bigl(2^{H(X_1)}+2^{H(X_2)}\bigr),
\]
即
\[
2^{H(X)}\le 2^{H(X_1)}+2^{H(X_2)}.
\]

\vspace{0.8em}

\begin{problem}{2.19}
设 $X$ 是 $[-1,1]$ 上的均匀分布的随机变量，试求 $H_c(X)$ 和 $H_c(X^2)$。
\end{problem}
\noindent\textbf{解：}
由于 $X$ 在 $[-1,1]$ 上均匀分布，
\[
f_X(x)=\frac12,\qquad -1\le x\le 1.
\]
所以
\[
H_c(X)=-\int_{-1}^{1}\frac12\log\frac12\,\mathrm{d}x=\log 2=1.
\]

令 $Y=X^2$。当 $0<y\le 1$ 时，
\[
f_Y(y)=f_X(\sqrt y)\frac{1}{2\sqrt y}
      +f_X(-\sqrt y)\frac{1}{2\sqrt y}
      =\frac{1}{2\sqrt y}.
\]
因此
\[
H_c(X^2)=H_c(Y)
=-\int_0^1\frac{1}{2\sqrt y}\log\frac{1}{2\sqrt y}\,\mathrm{d}y.
\]
令 $y=t^2$，则 $\mathrm{d}y=2t\,\mathrm{d}t$，得
\[
H_c(X^2)
=-\int_0^1\log\frac{1}{2t}\,\mathrm{d}t
=\int_0^1\log(2t)\,\mathrm{d}t
=1-\log e.
\]
所以
\[
H_c(X)=1,\qquad H_c(X^2)=1-\log e.
\]

\vspace{0.8em}

\begin{problem}{2.20}
令 $X$ 是取值 $\pm 1$ 的二元随机变量，概率分布为 $p(x=1)=p(x=-1)=0.5$，令 $Y$ 是连续随机变量，已知条件概率密度为
\[
p(y|x) =
\begin{cases}
\frac{1}{4} & -2 < y - x \leq 2 \\
0 & \text{其他}
\end{cases}
\]
试求：

\begin{enumerate}[label=\textnormal{(\alph*)}]
  \item $Y$ 的概率密度。

  \item $I(X;Y)$。

  \item 若对 $Y$ 做硬判决
  \[
  V =
  \begin{cases}
  1 & y > 1 \\
  0 & -1 < y \leq 1 \\
  -1 & y \leq -1
  \end{cases}
  \]
  求 $I(V;X)$，并对结果加以解释。
\end{enumerate}
\end{problem}
\noindent\textbf{解：}
\begin{enumerate}[label=\textnormal{(\alph*)}]
  \item 当 $X=1$ 时，$Y$ 在 $(-1,3]$ 上均匀分布；当 $X=-1$ 时，$Y$ 在 $(-3,1]$ 上均匀分布。因此
  \[
  p_Y(y)=\frac12p(y|x=1)+\frac12p(y|x=-1)
  =
  \begin{cases}
  \frac18, & -3<y\le -1,\\
  \frac14, & -1<y\le 1,\\
  \frac18, & 1<y\le 3,\\
  0, & \text{其他}.
  \end{cases}
  \]

  \item 由条件分布可知
  \[
  H_c(Y|X)=\log 4=2.
  \]
  又
  \[
  \begin{aligned}
  H_c(Y)
  &=-2\cdot\frac18\log\frac18
    -2\cdot\frac14\log\frac14
    -2\cdot\frac18\log\frac18\\
  &=\frac34+1+\frac34
  =\frac52.
  \end{aligned}
  \]
  故
  \[
  I(X;Y)=H_c(Y)-H_c(Y|X)=\frac52-2=\frac12.
  \]

  \item 由硬判决规则可得
  \[
  P(V=1|X=1)=\frac12,\quad P(V=0|X=1)=\frac12,\quad P(V=-1|X=1)=0,
  \]
  \[
  P(V=-1|X=-1)=\frac12,\quad P(V=0|X=-1)=\frac12,\quad P(V=1|X=-1)=0.
  \]
  所以
  \[
  P(V=1)=\frac14,\qquad P(V=0)=\frac12,\qquad P(V=-1)=\frac14.
  \]
  因此
  \[
  H(V)=H\!\left(\frac14,\frac12,\frac14\right)=\frac32,\qquad H(V|X)=1,
  \]
  从而
  \[
  I(V;X)=H(V)-H(V|X)=\frac12.
  \]
  这与 $I(X;Y)$ 相同。原因是区间 $(-1,1]$ 是两个条件密度的重叠区域，在该区域内 $Y$ 不能区分 $X$；而两侧非重叠区域已经能唯一判定 $X$。硬判决 $V$ 正好保留了这一区分信息，因此没有损失关于 $X$ 的互信息。
\end{enumerate}

\vspace{0.8em}

\begin{problem}{2.23}
一阶马尔可夫信源的状态图如习题2.23图所示，信源符号集为$\{0,1,2\}$，并定义$\bar{p}=1-p$。

\begin{enumerate}[label=\textnormal{(\alph*)}]
  \item 求信源平稳概率分布 $p(0), p(1), p(2)$。

  \item 求此信源的熵率 $H_\infty$。

  \item 近似认为此信源为无记忆时，符号概率分布等于平稳分布，求此近似信源的熵 $H(X)$，并与 $H_\infty$ 进行比较。

  \item 对一阶马尔可夫信源，$p$ 取何值时 $H_\infty$ 取最大值？又当 $p=0$ 和 $p=1$ 时结果如何？
\end{enumerate}
\end{problem}
\noindent\textbf{解：}
由状态图对应的转移关系可写出转移概率矩阵
\[
\boldsymbol{P}=
\begin{pmatrix}
1-p & \frac p2 & \frac p2\\
\frac p2 & 1-p & \frac p2\\
\frac p2 & \frac p2 & 1-p
\end{pmatrix}.
\]
\begin{enumerate}[label=\textnormal{(\alph*)}]
  \item 该矩阵各行、各列之和均为 $1$，且三个状态完全对称。因此当 $p>0$ 时，平稳分布唯一且为
  \[
  p(0)=p(1)=p(2)=\frac13.
  \]
  当 $p=0$ 时，$\boldsymbol{P}$ 为单位矩阵，平稳分布不唯一；若按对称平稳分布取值，仍有 $p(0)=p(1)=p(2)=1/3$。

  \item 一阶马尔可夫信源的熵率为
  \[
  H_\infty=\sum_i p(i)H(X_n|X_{n-1}=i).
  \]
  由于每一行的转移概率分布都是
  \[
  \left(1-p,\frac p2,\frac p2\right),
  \]
  故
  \[
  H_\infty
  =H\!\left(1-p,\frac p2,\frac p2\right)
  =-(1-p)\log(1-p)-p\log\frac p2.
  \]

  \item 若近似认为此信源无记忆，则符号概率分布为平稳分布，所以
  \[
  H(X)=-3\cdot\frac13\log\frac13=\log 3.
  \]
  因为
  \[
  H_\infty=H(X_n|X_{n-1})\le H(X_n)=H(X),
  \]
  所以
  \[
  H_\infty\le H(X)=\log 3.
  \]
  等号成立当且仅当下一状态与当前状态独立，即
  \[
  1-p=\frac p2=\frac13,
  \]
  因而 $p=2/3$。

  \item $H_\infty$ 是三元分布
  \[
  \left(1-p,\frac p2,\frac p2\right)
  \]
  的熵，当三项相等时取得最大值。因此
  \[
  1-p=\frac p2 \quad\Longrightarrow\quad p=\frac23,
  \]
  且
  \[
  H_{\infty,\max}=\log 3.
  \]
  当 $p=0$ 时，每个状态只转移到自身，故
  \[
  H_\infty=0.
  \]
  当 $p=1$ 时，每个状态以 $1/2$ 的概率转移到另外两个状态之一，故
  \[
  H_\infty=\log 2=1.
  \]
\end{enumerate}

\vspace{0.8em}

\begin{problem}{2.25}
\begin{enumerate}[label=\textnormal{(\alph*)}]
  \item 求状态转移矩阵为
  \[
  \boldsymbol{P} = \begin{pmatrix}
  1 - p_{01} & p_{01} \\
  p_{10} & 1 - p_{10}
  \end{pmatrix}
  \]
  的2状态马尔可夫信源的熵率。

  \item 使(a)中熵率最大的 $p_{01}, p_{10}$ 为多少？

  \item 如状态转移矩阵为
  \[
  \boldsymbol{P} = \begin{pmatrix}
  1 - p & p \\
  1 & 0
  \end{pmatrix}
  \]
  则相应的2状态马尔可夫信源熵率为多少？

  \item 寻找最大化(c)中熵率的 $p$ 值。

  \item 令 $N(t)$ 表示(c)中马尔可夫信源输出的长度为 $t$ 的可允许状态序列的数目。求 $N(t)$ 并计算 $H = \lim_{t \to \infty} \frac{1}{t} \log N(t)$。
\end{enumerate}
\end{problem}
\noindent\textbf{解：}
记二元熵函数为
\[
h(q)=-q\log q-(1-q)\log(1-q).
\]
\begin{enumerate}[label=\textnormal{(\alph*)}]
  \item 平稳分布满足
  \[
  \pi_0p_{01}=\pi_1p_{10},\qquad \pi_0+\pi_1=1.
  \]
  当 $p_{01}+p_{10}>0$ 时，
  \[
  \pi_0=\frac{p_{10}}{p_{01}+p_{10}},\qquad
  \pi_1=\frac{p_{01}}{p_{01}+p_{10}}.
  \]
  因此熵率为
  \[
  H_\infty
  =\pi_0 h(p_{01})+\pi_1 h(p_{10})
  =\frac{p_{10}h(p_{01})+p_{01}h(p_{10})}{p_{01}+p_{10}}.
  \]
  当 $p_{01}=p_{10}=0$ 时，两状态均为吸收状态，熵率为 $0$。

  \item 因为每一行的条件熵都不超过 $1$，所以
  \[
  H_\infty\le 1.
  \]
  当且仅当两行转移概率都为均匀分布时取到最大值，即
  \[
  p_{01}=p_{10}=\frac12,
  \]
  此时 $H_{\infty,\max}=1$。

  \item 由
  \[
  \boldsymbol{P} =
  \begin{pmatrix}
  1-p & p\\
  1 & 0
  \end{pmatrix}
  \]
  可知平稳分布满足
  \[
  \pi_1=\pi_0p,\qquad \pi_0+\pi_1=1,
  \]
  因而
  \[
  \pi_0=\frac{1}{1+p},\qquad \pi_1=\frac{p}{1+p}.
  \]
  状态 $0$ 下的条件熵为 $h(p)$，状态 $1$ 下下一状态确定为 $0$，条件熵为 $0$。故
  \[
  H_\infty=\frac{h(p)}{1+p}.
  \]

  \item 对
  \[
  H_\infty(p)=\frac{h(p)}{1+p}
  \]
  求极大值。令导数为 $0$，得
  \[
  (1+p)\log\frac{1-p}{p}=h(p).
  \]
  化简为
  \[
  (1-p)^2=p.
  \]
  解得
  \[
  p=\frac{3-\sqrt5}{2}.
  \]
  因为端点 $p=0,1$ 时熵率均为 $0$，故该点为最大值点。

  \item 由转移矩阵可知，状态 $1$ 后面不能再接状态 $1$。因此 $N(t)$ 等于长度为 $t$ 且不含连续两个 $1$ 的二元序列数。

  设 $N(1)=2,N(2)=3$。对 $t\ge 3$，若序列以 $0$ 结尾，则前 $t-1$ 位任意可允许；若以 $1$ 结尾，则第 $t-1$ 位必须为 $0$，前 $t-2$ 位任意可允许。因此
  \[
  N(t)=N(t-1)+N(t-2).
  \]
  所以
  \[
  N(t)=F_{t+2},
  \]
  其中 $F_0=0,F_1=1$。于是
  \[
  H=\lim_{t\to\infty}\frac1t\log N(t)
  =\log\frac{1+\sqrt5}{2}.
  \]
\end{enumerate}

\end{document}
